\documentclass[10pt,a4paper]{amsart}
\usepackage[margin=19mm]{geometry}
\usepackage{amsmath,amssymb,mathtools}
\usepackage[scr=rsfs,cal=euler]{mathalfa}
\usepackage{xfrac}
\usepackage[unicode,hidelinks,bookmarks=false]{hyperref}
\newcommand{\ie}{i.e.\ }
\newcommand{\cf}{cf.\ }
\newcommand{\resp}{respectively\ }
\newcommand{\Subsection}{\subsection}
\providecommand{\nobreakdash}{\nobreak\hbox{-}}
\setlength{\emergencystretch}{2em}
\multlinegap0pt
\usepackage{amssymb}
\usepackage{mathtools}
\usepackage[shortlabels]{enumitem}
\usepackage{tikz}
\usepackage{float}
\usepackage{xcolor}
\newtheorem{restatable}{Restatable}
\newtheorem{proposition}{Proposition}
\newtheorem{claim}{Claim}
\DeclareMathOperator{\id}{id}
\newcommand{\idem}{\mathrm{idem}}
\newcommand{\pol}{\operatorname{Pol}}
\newcommand{\str}[1]{\mathbb{#1}}
\title{Cut-homotopies and the complexity\\ of edge-coloring problems}
\author{Alexey Barsukov, Roman Feller, Maximilian Hadek, Davide Perinti}
\date{Source: arXiv:2607.09631v1 (CC BY 4.0)}
\begin{document}
\maketitle

\noindent\emph{Source and licence.} Cut-homotopies and the complexity\\ of edge-coloring problems. Version arXiv:2607.09631v1 (CC BY 4.0), distributed by arXiv under the Creative Commons Attribution 4.0 International licence, http://creativecommons.org/licenses/by/4.0/, as recorded in the arXiv licence metadata for this exact version and shown on the abstract page https://arxiv.org/abs/2607.09631v1. Full licence notice: \texttt{licenses/arxiv-2607.09631-CC-BY-4.0.txt}.\par\smallskip

\noindent\textit{Editorial scope.} Counted unit: the article's proposition prop:noHomotopies (source0.tex lines 1955--1958). The article's complete original proof of it is delivered with it (source0.tex lines 1960--2044, one \texttt{proof} environment of 85 lines). No mathematics is written, repaired or re-derived: the statement and the proof are the article's own lines, in the article's own notation, and no retained line is substituted or re-flowed.  Retained uncounted as the source-local context the proof needs: lines 1942--1951.  Nothing was omitted from any retained span, so the package carries no omission record.  No retained line contains a citation, so the wrapper carries no bibliography and no bibliography entry.  No retained line contains a figure environment, an image-inclusion command or drawing code, so no image is delivered and the package declares no asset dependency.  Editorial formatting only, in addition: an AMS \texttt{amsart} preamble that restates the article's own declarations and macros used by the retained lines, namely the environments \texttt{restatable}, \texttt{proposition}, \texttt{claim} on separate counters, together with the standard packages those lines need.  The retained environments print in this document as Restatable 1, Proposition 1, Claim 1 in order of appearance, whereas the article prints the same items with the numbering of its own full document.  The complete native source of the article as distributed in the arXiv e-print for this exact version is bundled unchanged as \texttt{sources/204884/source0.tex}.
\begin{restatable}{lemma}{mapFactory}
\label{lem:mapFactory}
    Let $n$ be a number, and for each $a\in \str A^n$ pick a subset $U_{a}\subseteq A$. Assume that for every graph $G$ and every coloring $\chi\colon \hat G\to \str A^n$, there is a coloring $\eta\colon \hat G\to \str A$ such that 
    \begin{equation*}
        \forall vw\in\hat G: \eta(vw)\in U_{\chi(vw)}.
    \end{equation*}
    Then there exists a polymorphism 
    $f \colon \str A^n \to \str A$ 
    with $f(a)\in U_{a}$ for all $a \in \str A^n$.
\end{restatable}

\begin{proposition}
\label{prop:noHomotopies}
    Let $\str A$ be a core and $(S,\theta)$ be a naked set of $\pol(\str A)_\idem$. Then the minion homomorphism $\alpha\colon \pol(\str A)_\idem\to\id$ induced by $(S, \theta)$ collapses idempotent cut-homotopies, i.e., if $f,g\in\pol(\str A)_\idem$ and $f\sim_1 g$ with $h$ idempotent, then $\alpha(f)=\alpha(g)$.
\end{proposition}

\begin{proof}

Striving for a contradiction we assume that $h$ is an idempotent cut-homotopy between $n$-ary polymorphisms $h_0,h_1$ that act as different projections on $S/_\theta$, without loss of generality, the first and second projection. Then, $h' \colon \str A^2 \times \hat P_1 \to \str A, (x,y,ij) \mapsto h_{ij}(x,y,\dots, y)$ is a cut-homotopy between binary polymorphisms of $\str A$ that act as different projections on $S/_\theta$. Hence, we may assume $h$ to be binary, and that $h_0, h_1$ act on $S/_\theta$ as the first and second projection, respectively.



Let $d$ be the homomorphism $(h_0, h_1) \colon \str A^2\to \str A^2, (a,b) \mapsto (h_0(a,b),h_1(a,b))$, and choose $N$ so that $d^N$ is a retraction, i.e., $d^N = d^{2N}$. Let $D$ be the image of $d^N$, and note that $d^N$ restricts to the identity on $D$.

\begin{claim}
The composition $f\coloneq h_{01}\circ d^N$ is an idempotent polymorphism of $\str A$ satisfying $f|_D = h_{01}|_D$.
\end{claim}
\begin{proof}[Proof of the claim]
We show that for every $K_m$ and every coloring $\chi\colon \hat K_m\to \str A^2$, 
the map $f \circ \chi$ is a homomorphism $\hat K_m \to \str A$.
Define $\chi_0\coloneq d^N\circ \chi$.

Fix a vertex $v\in K_m$ and consider the two cuts $c_{0} \coloneq K_m\setminus\{v\}\mid v$ and
$c_1\coloneq  v\mid K_m\setminus\{v\}$.
Applying $h$ to $\chi_0$ along each cut yields a coloring $\hat K_m \to \str A$ and we let $\chi_1\colon \hat K_m \to \str A^2$ denote their product. More formally, $\chi_1$ is defined to be the homomorphism $\big( h \circ (\chi_0,\hat c_0), h \circ (\chi_0,\hat c_1) \big)\colon \hat K_m \to \str A^2$.
For every $w w' \in \hat K_m$ it holds
    \begin{equation*}
            \chi_1(w w') =
    \begin{cases}
    \big(h_0\circ\chi_0(w w'), h_1\circ \chi_0(w w') \big) = d\circ \chi_0(w w') &\text{if } v\notin \{w, w'\} \\ 
    \big(h_{01}\circ\chi_0(w w'), h_{01}\circ\chi_0(w w')\big) &\text{if } v\in \{w, w'\}.
    \end{cases}
    \end{equation*}
    Observe that in the latter case both entries of $\chi_1(w w')$ agree. We iterate this construction, obtaining $\chi_{i+1}$ from $\chi_i$ in the same way $\chi_1$ is obtained from $\chi_0$ until we arrive at the coloring $\chi_N$. The aforementioned observation in combination with the fact that $h$ is idempotent guarantees that for every $w w' \in \hat K_m$ it holds 
    \begin{equation*}
            \chi_N(w w') =
    \begin{cases}
    d^N \circ \chi_0(w w') = \chi_0(w w')&\text{if } v\notin \{w, w'\}\\ 
    \big(h_{01}\circ\chi_0(w w'),h_{01}\circ\chi_0(w w')\big) &\text{if } v\in \{w, w'\}.
    \end{cases} 
    \end{equation*}
    Replacing $\chi_0$ with $\chi_N$ we repeat this procedure for every vertex of 
    $K_m$.
    Once we have exhausted all vertices, we call the resulting coloring $\chi'\colon \hat K_m \to \str A^2$. For every edge $w w'$ of $K_m$, as it is incident to precisely two vertices of $K_m$, we have
    \begin{align*}
        \chi'(w w') &= \Big( h_{01}\big(h_{01}\circ \chi_0(w w'), h_{01}\circ \chi_0(w w')\big), h_{01}\big(h_{01}\circ \chi_0(w w'), h_{01}\circ \chi_0(w w')\big) \Big)\\
        &= \big( h_{01}\circ \chi_0(w w'), h_{01}\circ \chi_0(w w') \big)\\
        &= \big( f \circ \chi(w w'), f \circ \chi(w w') \big),
    \end{align*}
    in particular, $f \circ \chi$ is a homomorphism $\hat K_m \to \str A$.
    Thus, $f$ is a homomorphism; its idempotence is clear.
\end{proof}

  
Without loss of generality we may assume that $f$ acts as the first projection on $S/_\theta$. Thus, by definition of $f$ each pair $(a,b) \in S^2 \cap D$ satisfies $h_{01}(a,b)/_\theta = a/_\theta$.
Further, recall that $h_0$ and $h_1$ act as the first and second projection on $S/_\theta$, respectively.
 
We now fix $a,b\in S$ in different classes of $\theta$.
We use Lemma~\ref{lem:mapFactory} to show that there is a polymorphism $g: \str A^2 \rightarrow \str A$ with $g(a,b)/_\theta = b/_\theta$ and $g(a',b')/_\theta = a'/_\theta$ whenever $a',b'\in S$ with $(a,b)/_\theta\not= (a',b')/_\theta$, contradicting the fact that $(S,\theta)$ is a naked set.
To this end, we verify that for every graph $G$ and every coloring $\chi\colon \hat G\to \str A^2$, there is a coloring $\eta\colon \hat G\to \str A$ such that for all $vw \in \hat G$ it holds:
\begin{equation*}
    \eta(vw) \in 
\begin{cases}
    b/_\theta, &\text{ if } \chi(v w) \in S^2 \text{ and } \chi(v w)/_\theta = (a,b)/_\theta \\
    a'/_\theta, &\text{ if } \chi(vw) = (a',b') \in S^2\text{ and } (a',b')/_\theta \not= (a,b)/_\theta.
\end{cases}
\end{equation*}

So let $G$ be a graph and $\chi = (\chi_1,\chi_2) \colon \hat G\to \str A^2$ a coloring. Note that the coloring $d^N\circ \chi$ agrees with $\chi$ modulo
$(S,\theta)$ on all $vw \in \hat G$ with $\chi(vw)\in S^2$ and additionally its image is contained in $D$. 
Take any edge $v_0 w_0 \in \hat G$ with $\chi(v_0 w_0)/_\theta = (a,b)/_\theta$ and consider the cut $c \coloneq G\setminus \{v_0,w_0\} \mid \{v_0,w_0\}$. 
Define $\eta'\colon \hat G\to \str A$ to be the coloring $h \circ (d^N\circ \chi,\hat c)$. Using the properties of $h$, in particular that $h_{01}|_D$ acts as the first projection on $S/_\theta$ we obtain that for $v w \in \hat G$ it holds
\begin{equation*}
    \eta'(v w)  \in
\begin{cases}
    h_{1}(d^N \circ \chi(v w))/_\theta = b/_\theta &\text{ if } v w = v_0 w_0 \\
    h_{01}(d^N \circ \chi(v w))/_\theta =  a'/_\theta &\text{ if } \chi(v w) = (a',b') \in S^2\text{ and } |\{v, w\}\cap \{v_0,w_0\}| = 1 \\
    h_{0}(d^N \circ \chi(v w))/_\theta = a'/_\theta &\text{ if } \chi(v w) = (a',b') \in S^2 \text{ and } |\{v, w\} \cap \{v_0,w_0\}| = 0.
\end{cases}
\end{equation*}
Define the coloring $\chi' \coloneq (\eta',\chi_2)\colon \hat G \to \str A$ and observe that on all $vw \in \hat G\setminus \{v_0 w_0\}$ with $\chi(vw)\in S^2$, $\chi(vw)$ and $\chi'(vw)$ agree modulo $\theta$ and furthermore $\chi'(v_0 w_0)/_\theta = (b,b)/_\theta$. 
We repeat the same procedure with $\chi'$ in place of $\chi$ with every further edge $v_0' w_0' \in \hat G$ with $\chi(v_0' w_0')/_\theta = (a,b)/_\theta$ until we have exhausted all such edges. We obtain a coloring $\tilde \eta \colon \hat G \to \str A$ such that for all $vw \in \hat G$ it holds
\begin{equation*}
        \tilde \eta(v w) \in 
\begin{cases}
    b/_\theta &\text{ if } \chi(v w)/_\theta = (a,b)/_\theta\\
    a'/_\theta &\text{ if } \chi(v w) = (a',b') \in S^2 \text{ and } (a',b')/_\theta \not= (a, b)/_\theta.
\end{cases}
\end{equation*}
Thus, $\tilde \eta$ is the desired coloring $\eta$.
\end{proof}

\end{document}
